\section{Corrected Problem Formulation}
\label{sec:setup}
Let \(A=[a_1,\ldots,a_n]\in\R^{p\times n}\), let \(X=[x_1,\ldots,x_m]\in\St(p,m)\), and let \(\mu_j>0\). For a vector \(c\), define soft thresholding by
\[
\Soft_\gamma(c)_i=\operatorname{sign}(c_i)(|c_i|-\gamma)_+.
\]
The component response is \(c_j(X)=\mu_jA^\top x_j\).

\subsection{A-SUM: normalized loading-ball model}
The corrected primal is
\begin{equation}
\max_{X\in\St(p,m)}\sum_{j=1}^m
\max_{\norm{z_j}_2\le 1}
\bigl\{c_j(X)^\top z_j-\gamma_j\norm{z_j}_1\bigr\}.
\label{eq:asum-primal}
\end{equation}
Loading balls allow an inactive column. The reduced objective is
\begin{equation}
F_{\mathrm{sum}}(X)=\sum_{j=1}^m\norm{\Soft_{\gamma_j}(c_j(X))}_2.
\label{eq:asum}
\end{equation}
When the thresholded response is nonzero, the recovered loading is its normalized version; otherwise \(z_j=0\) is valid.

\begin{proposition}[Exact loading-ball elimination]
\label{prop:sum-elimination}
For any \(c\in\R^n\) and \(\gamma\ge0\),
\[
\max_{\norm z_2\le1}\{c^\top z-\gamma\norm z_1\}
=\norm{\Soft_\gamma(c)}_2.
\]
If \(\Soft_\gamma(c)\neq0\), the unique maximizer is \(z^\star=\Soft_\gamma(c)/\norm{\Soft_\gamma(c)}_2\); if it is zero, every maximizer includes \(z=0\), and zero is selected by convention.
\end{proposition}
\begin{proof}
For each coordinate, \(c_i z_i-\gamma|z_i|\le (|c_i|-\gamma)_+|z_i|\), with equality when the sign of \(z_i\) agrees with \(c_i\) on the active set and \(z_i=0\) elsewhere. Hence the inner problem reduces to maximizing \(\Soft_\gamma(c)^\top u\) over \(\norm u_2\le1\). Cauchy--Schwarz yields the value and optimizer. If the thresholded vector is zero, the objective is nonpositive and zero attains zero.
\end{proof}

\subsection{A-SQ: squared-threshold model}
The historical reduced objective is instead
\begin{equation}
F_{\mathrm{sq}}(X)=\sum_{j=1}^m\norm{\Soft_{\gamma_j}(c_j(X))}_2^2.
\label{eq:asq}
\end{equation}
It has a different primal interpretation.
\begin{proposition}[Quadratic-lift identity]
\label{prop:quadratic-lift}
For any \(c\in\R^n\),
\[
\norm{\Soft_\gamma(c)}_2^2
=\max_{z\in\R^n}\{2c^\top z-\norm z_2^2-2\gamma\norm z_1\},
\]
and the unique maximizer is \(z^\star=\Soft_\gamma(c)\).
\end{proposition}
\begin{proof}
The objective separates over coordinates. For one coordinate, completing the square on each sign half-line gives the maximizer \(\operatorname{sign}(c_i)(|c_i|-\gamma)_+\) and optimal value \((|c_i|-\gamma)_+^2\). Summing proves the claim.
\end{proof}

\subsection{The objectives are not equivalent}
\begin{proposition}[Two-component counterexample]
\label{prop:counterexample}
There exist a compact feasible set and two nonnegative component responses \(r_1(X),r_2(X)\) for which \(\arg\max_X r_1+r_2\) and \(\arg\max_X r_1^2+r_2^2\) are disjoint.
\end{proposition}
\begin{proof}
Consider two feasible points \(X_a,X_b\) with responses \((r_1,r_2)(X_a)=(2,2)\) and \((r_1,r_2)(X_b)=(3.2,0.5)\). The sums are \(4\) and \(3.7\), so the sum objective selects \(X_a\). The squared sums are \(8\) and \(10.49\), so the squared objective selects \(X_b\). Such response pairs can be realized by a finite feasible subset or approximated arbitrarily closely by thresholded linear responses. Therefore no monotone scalar transformation relates the block objectives in general.
\end{proof}

% §14 required table #1, models and loading semantics.
\begin{table}[t]
\centering
\caption{Corrected block estimator family and loading recovery. The three columns define \emph{distinct} estimators; §3 proves A-SUM $\neq$ A-SQ at $m\!\ge\!2$ and A-HT is a third exact-cardinality variant. Loading recovery is the loading-column argmax at a fixed score and defines the \emph{raw} estimand that the restricted refit stage then post-processes.}
\label{tab:a-models-and-loadings}
\small
\begin{tabular}{p{0.17\linewidth}p{0.26\linewidth}p{0.23\linewidth}p{0.24\linewidth}}
\toprule
Estimator & Primal (per column $j$) & Reduced objective & Loading recovery / inactive policy \\
\midrule
A-SUM &
$\displaystyle \max_{\norm{z_j}_2\le 1}\!\!\bigl\{c_j^{\!\top}z_j-\gamma_j\norm{z_j}_1\bigr\}$ &
$F_{\mathrm{sum}}(X)=\sum_j\norm{\Soft_{\gamma_j}(c_j)}_2$ &
$z_j^\star = \Soft_{\gamma_j}(c_j)/\norm{\Soft_{\gamma_j}(c_j)}_2$ when nonzero; otherwise $z_j=0$ is the chosen inactive branch (loading-ball slack) \\
\addlinespace
A-SQ &
$\displaystyle \max_{z_j}\!\!\bigl\{2c_j^{\!\top}z_j-\norm{z_j}_2^2-2\gamma_j\norm{z_j}_1\bigr\}$ &
$F_{\mathrm{sq}}(X)=\sum_j\norm{\Soft_{\gamma_j}(c_j)}_2^2$ &
$z_j^\star = \Soft_{\gamma_j}(c_j)$ (unconstrained quadratic lift; the vector itself is the loading, no renormalization) \\
\addlinespace
A-HT (hard-threshold, exact cardinality~$k$) &
$\displaystyle \max_{\norm{z_j}_2\le1,\,\norm{z_j}_0\le k}\!\!c_j^{\!\top}z_j$ &
$F_k(X)=\sum_j\sum_{i\in T_k(c_j)}c_{j,i}^2$ &
$z_j^\star$ supported on the top-$k$ coordinates of $c_j$ (deterministic tie policy); refit on the chosen support \\
\bottomrule
\end{tabular}
\end{table}


\subsection{Numerical conventions}
The algorithmic object is incomplete without deterministic edge-case policies. We use: (i) inactive response columns remain zero and are excluded from the active polar block; (ii) threshold equality is inactive unless an experiment explicitly studies set-valued ties; (iii) rank-deficient polar factors are chosen by a deterministic thin-SVD completion; (iv) component permutations are aligned by maximum absolute score correlation for diagnostics; and (v) refitting applies the support mask before normalization or solves the exact restricted problem.

% §14 required table #3, edge-case policies.
% Operational conventions the implementation (facet.*) and the proofs
% both commit to. A theorem is said to apply to the solver only when it
% is run under the exact policy named in its statement.
\begin{table}[t]
\centering
\caption{Edge-case policies in the FACET reference implementation. These are the \emph{operational} conventions on which every proved theorem is explicitly conditional; running the solver under a different policy voids the theorem's claim.}
\label{tab:a-edge-case-policies}
\small
\begin{tabular}{p{0.26\linewidth}p{0.32\linewidth}p{0.34\linewidth}}
\toprule
Edge case & Policy & Pinning test \\
\midrule
Threshold tie ($|c_i|=\gamma$) & Treated as inactive ($\Soft_\gamma(c)_i=0$) via the standard $(\cdot)_+$ convention; signed subgradient reused by the solver. & \texttt{tests/test\_paper\_a\_identities.py::test\_a1\_*} \\
\addlinespace
Zero soft-threshold vector ($\Soft_\gamma(c)=0$) & Loading set to zero; the loading-ball slack delivers primal value $0$; no renormalization attempted. & \texttt{test\_a1\_sum\_norm\_zero\_support\_returns\_zero} \\
\addlinespace
Mask-first refit on a non-empty support & $z = M c / \norm{Mc}_2$; empty support raises \texttt{ZeroDirectionError} rather than returning NaN. & \texttt{tests/test\_paper\_a\_edge\_cases.py::test\_a4\_mask\_then\_normalize\_*} \\
\addlinespace
Rank-deficient polar input (Stiefel oracle) & \texttt{canonical} policy completes via SVD; \texttt{raise} policy rejects and surfaces the assumption violation. Both treated as explicit theorem conditioning. & \texttt{test\_a5\_canonical\_policy\_*}, \texttt{test\_a5\_raise\_policy\_*} \\
\addlinespace
Zero direction to a linear oracle & Both polar policies raise \texttt{ZeroDirectionError} unconditionally. & \texttt{test\_a5\_zero\_direction\_raises\_regardless\_of\_policy} \\
\addlinespace
Top-$k$ tie at the boundary (A-HT) & Lower original index wins; set-valued tie metadata recorded on \texttt{TopKSelection}. & \texttt{test\_topk\_tie\_policy\_prefers\_lower\_index} \\
\addlinespace
Armijo backtracking at machine precision & Line search fails loudly (\texttt{RuntimeError}); drivers may catch and drop that seed rather than silently reporting a stationary point. & empirical in \texttt{tests/test\_paper\_a\_solver\_dynamics.py} \\
\addlinespace
Component sign at convergence & Component orientation is global-sign ambiguous; projector-level reporting is sign-invariant. & \texttt{test\_a9\_sign\_identification\_matches\_planted\_sign\_pattern} \\
\bottomrule
\end{tabular}
\end{table}

\begin{table}[t]
\centering
\caption{Assumption and failure-mode taxonomy.}
\label{tab:a-assumptions}
\small
\begin{tabular}{p{0.24\linewidth}p{0.34\linewidth}p{0.32\linewidth}}
\toprule
Assumption & Enables & Failure manifestation \\
\midrule
Strict threshold margin & Finite support/sign identification & Persistent switching or tie dependence \\
Full-rank local polar argument & Smooth polar branch & Set-valued update and unstable derivative \\
Active eigengap or $\rho(D\Phi)<1$ & Local linear convergence & Slow/noncontractive local dynamics \\
Certified distance radius & Safe screening & Heuristic false elimination if radius is guessed \\
Independent document units & Valid uncertainty summaries & Token-level pseudo-replication \\
\bottomrule
\end{tabular}
\end{table}

